\section{Metodika hodnocení}
Hra bude rozdělena na kola, přičemž každá hra se bude skládat z neurčitého počtu kol. Hra je navržena tak, aby hráč bojoval se svými vojáky proti nepřátelským vojákům, kteří jsou řízeni autonomním agentem. Když hra začne, hráč si postaví vlastní vojáky, kteří mu budou pomáhat v boji, a při vítězství nad nepřátelskými vojáky, jež budou vytvořeni podle předem naučeného SARSA algoritmu. 

V průběhu kola se také předem naučená AI část vojáků, vytvořená za pomoci Deep Q-Learning algoritmu, bude v průběhu bitvy upravovat (učit), aby se pokusila přelstít hráče a jeho vojáky a tím je porazit.

Na konci kola hráč využije sesbírané informace k rozhodnutí, zda chce nebo potřebuje upravit své vojáky pro další kolo. Agent využije své nasbírané informace k úpravě SARSA algoritmu, aby dokázal vybrat efektivnější výbavu proti hráči. Nakonec se tyto informace pošlou do Deep Q-Learning algoritmu a tím se upraví chování vojáků, aby byli schopni lépe a efektivněji porazit hráče a jeho vojáky.

Pokud hra neskončí, další kolo začne s těmito novými úpravami a novou výbavou. Tento cyklus se opakuje, dokud nenastane konec hry. Když ke konci hry dojde, algoritmy se revertují do svého základního, předem naučeného stavu. Diagram logiky celé hry lze vidět na obrázku \ref{obrDiaLog}.

Díky tomu, že v Deep Q-Learning algoritmu se používá hluboká neuronová síť (DNN), lze vyzkoušet různé neuronové sítě a zhodnotit, která z nich produkuje nejlepší výsledky a nejlepší herní zážitek. Po zjištění nejlepší DNN pro tento Deep Q-Learning algoritmus by bylo možné zhodnotit, který z typů hodnocení, popsaných v sekci \ref{TH}, by byl nejlepší pro herní zážitek. Tím by vznikla nejlepší varianta agenta, proti kterému by bylo zábavné bojovat.

\newpage

\begin{figure}[H]
\includegraphics[scale=0.75]{4-1.png}
\centering
\caption{Diagram logiky hry.}
\label{obrDiaLog}
\end{figure}

\section{Typy hodnocení}
\makeatletter\def\@currentlabel{4.2	Typy hodnocení}\makeatother
\label{TH}
Agenta lze hodnotit různými způsoby, ale byly zvoleny následující tři návrhy, protože budou mít nejvýraznější a nejzajímavější vliv na chování agenta umělé inteligence:

\subsection{Podle efektivity zabití nepřátel}
Pokud bude agent hodnocen podle množství nepřátel, které dokáže porazit, pravděpodobně si přizpůsobí své chování tak, že bude neustále útočit. Jeho vojáci se pravděpodobně zaměří na rychlost a schopnost rychle porazit nepřítele, i když budou mít velmi málo života a například horší vidění. Tito vojáci se zřejmě budou méně schovávat, protože by to snižovalo jejich rychlost a počet poražených nepřátel. Díky těmto změnám by agent mohl nasadit více vojáků, což by ztížilo boj hráči, který by se musel soustředit na více nepřátel zároveň.

Alternativně by agent mohl zvolit jiný přístup a vytvořit jednoho super vojáka, který by měl hodně života, silné brnění a zbraň s vysokým poškozením, ale byl by téměř nepohyblivý a měl ještě horší vidění. Tento voják by byl velmi těžko zničitelný, ale pokud by se ho podařilo zničit, agent by už nemohl nasadit další vojáky.

\subsection{Podle efektivity přežití}
Agent se pravděpodobně bude chovat zcela opačně, pokud bude hodnocen podle toho, jak dlouho dokáže jeden voják přežít, namísto toho, kolik nepřátelských vojáků porazí. V takovém případě by se soustředil na množství života a brnění, přičemž by ho nezajímalo, kolik poškození dokážou jeho vojáci udělit hráči nebo nepřátelským vojákům. To je jeden z možných směrů, kterými by se mohl vydat. Alternativně by mohl vytvořit vojáka, který se pohybuje rychle a útočí z dálky, například pomocí odstřelovací pušky, místo aby bojoval v blízkosti nepřátel. Agent se pravděpodobně nebude snažit nasadit ani malé, ani velké množství vojáků, aby maximalizoval šance na přežití těch ostatních.

\subsection{Kombinace předchozích hodnocení}
Nejzajímavější situace nastává, když se hodnotí jak přežití, tak efektivita zničení nepřítele. V této situaci je možné vytvořit mnoho různých scénářů, ale zároveň může být tato situace méně zajímavá, protože existuje vysoká šance, že agent najde scénář, který bude ideální pro toto hodnocení, a to mnohem rychleji než v obou předchozích případech. Například by agent mohl vytvořit \uv{super odstřelovače}, který by měl slušné množství života, brnění a rychlosti, a zároveň by se držel mimo dosah nepřátel a ničil je z dálky pomocí odstřelovací pušky. I když by byl zcela sám, pokud by hráč ani jeho tým neměli možnost tohoto vojáka z dálky zneškodnit, vedlo by to k velmi rychlé porážce hráče.